% ------------------------------------------------------------------
%  Preambolo delle dispense (skill dispense-latex, layout del 22/10/2026)
%  Modello: frontespizio + indice a pagina intera, etichette colorate
%  numerate per sezione, dimostrazioni a passi con barra laterale,
%  riquadro verde solo per i risultati essenziali.
%  Prima di % ------------------------------------------------------------------
%  Preambolo delle dispense (skill dispense-latex, layout del 22/10/2026)
%  Modello: frontespizio + indice a pagina intera, etichette colorate
%  numerate per sezione, dimostrazioni a passi con barra laterale,
%  riquadro verde solo per i risultati essenziali.
%  Prima di % ------------------------------------------------------------------
%  Preambolo delle dispense (skill dispense-latex, layout del 22/10/2026)
%  Modello: frontespizio + indice a pagina intera, etichette colorate
%  numerate per sezione, dimostrazioni a passi con barra laterale,
%  riquadro verde solo per i risultati essenziali.
%  Prima di % ------------------------------------------------------------------
%  Preambolo delle dispense (skill dispense-latex, layout del 22/10/2026)
%  Modello: frontespizio + indice a pagina intera, etichette colorate
%  numerate per sezione, dimostrazioni a passi con barra laterale,
%  riquadro verde solo per i risultati essenziali.
%  Prima di \input{preambolo} definisci \corso, \argomento, \lezioni.
%  Compilare con pdflatex (scripts/setup_pdflatex_it.sh per la sillabazione).
% ------------------------------------------------------------------
\providecommand{\corso}{Nome del corso}
\providecommand{\argomento}{Argomento}
\providecommand{\lezioni}{Lezioni}
\usepackage[T1]{fontenc}
\usepackage{babel}
\babelprovide[import, main]{italian}
\usepackage{lmodern}            % Latin Modern (se manca: apt-get install -y lmodern)
\usepackage{microtype}
\usepackage[a4paper, left=2.4cm, right=2.4cm, top=2.5cm, bottom=2.7cm]{geometry}
\usepackage[skip=6pt plus 2pt, indent=0pt]{parskip}

\usepackage{amsmath, amssymb}
\usepackage{bm}
\numberwithin{equation}{section}
\usepackage{siunitx}
\sisetup{output-decimal-marker={,}, group-separator={\,}, exponent-product=\cdot}
\DeclareSIUnit{\angstrom}{\text{\AA}}
\usepackage{booktabs, array, tabularx, multirow, longtable}
\usepackage{enumitem}
\setlist{itemsep=3pt, topsep=4pt, parsep=0pt}
\usepackage[font=small, labelfont=bf, format=hang, skip=6pt]{caption}
\usepackage{float}
\usepackage{needspace}
% \spazio{n\baselineskip}: va a capo pagina solo se restano meno di n righe.
% Sostituisce \needspace, il cui \penalty-100 con \vskip elastico fa spezzare
% le pagine prima dei titoli anche quando c'è posto (pagine mezze vuote).
\makeatletter
\newcommand{\spazio}[1]{\par\begingroup\setlength{\dimen@}{#1}%
  \dimen@ii\pagegoal\advance\dimen@ii-\pagetotal
  \ifdim\dimen@>\dimen@ii\ifdim\dimen@ii>\z@\vfil\fi\penalty-10000\fi\endgroup}
\makeatother

% ---------------------------------------------------------- colori
\usepackage{xcolor}
\definecolor{csez}{HTML}{193C82}   % titoli di sezione
\definecolor{csub}{HTML}{3C5995}   % sottosezioni, link
\definecolor{cdef}{HTML}{1E5AAA}   % Definizione
\definecolor{cris}{HTML}{14823C}   % Risultato, riquadro verde
\definecolor{cbox}{HTML}{F1F7F3}   % fondo del riquadro verde
\definecolor{cform}{HTML}{E4F0E8}  % fondo di \formula
\definecolor{cformb}{HTML}{ACD3BA} % bordo di \formula
\definecolor{coss}{HTML}{BE1E1E}   % Osservazione (refusi, convenzioni)
\definecolor{cint}{HTML}{007D87}   % Interpretazione fisica, «Cosa dice»
\definecolor{cappr}{HTML}{783296}  % Approfondimento
\definecolor{crie}{HTML}{C86E00}   % Riepilogo
\definecolor{cdim}{HTML}{6E6E6E}   % Dimostrazione
\definecolor{cpasso}{HTML}{585858} % Passo n.
\definecolor{cesm}{HTML}{4A6FA5}   % Esempio
\definecolor{cenb}{HTML}{EEF3FA}   % fondo del riquadro blu dell'enunciato
% colori per le figure
\definecolor{accent}{HTML}{193C82}
\definecolor{vA}{HTML}{C0392B}
\definecolor{vB}{HTML}{2E6FB7}
\definecolor{vC}{HTML}{23884A}
\definecolor{atomA}{HTML}{E15759}
\definecolor{atomB}{HTML}{4E79A7}
\definecolor{atomC}{HTML}{F28E2B}
\definecolor{atomD}{HTML}{59A14F}
\definecolor{atomE}{HTML}{B07AA1}
\definecolor{cellP}{HTML}{3C9D4E}
\definecolor{cellC}{HTML}{1F4E8C}
\definecolor{cellS}{HTML}{7B3FA0}
\definecolor{bad}{HTML}{C0392B}
\definecolor{Na}{HTML}{F28E2B}
\definecolor{Cl}{HTML}{4E79A7}
% nomi del preambolo precedente (figure e snippet esistenti)
\colorlet{boxdef}{cdef}\colorlet{boxkey}{cris}\colorlet{boxwarn}{coss}
\colorlet{boxex}{black!65}\colorlet{boxdeep}{cappr}\colorlet{mapgray}{black!60}

% ---------------------------------------------------------- figure
\usepackage{tikz}
\usetikzlibrary{arrows.meta, calc, decorations.pathreplacing, patterns}
\usepackage{pgfplots}
\pgfplotsset{compat=1.18, every axis/.append style={font=\footnotesize}, trig format plots=rad}
\usepackage[europeanresistors, cuteinductors, europeanvoltages, straightvoltages]{circuitikz}
\ctikzset{bipoles/length=1.0cm}
\usepackage{tikz-3dplot}
\tikzset{
  vec/.style={-{Stealth[length=2.4mm, width=1.7mm]}, thick},
  bond/.style={black!22, line width=0.5pt},
  latline/.style={black!30, thin},
  edge/.style={black!75, line width=0.7pt},
  edgehid/.style={black!55, line width=0.6pt, densely dashed},
  edgelight/.style={black!50, line width=0.6pt},
  edgelighthid/.style={black!40, line width=0.5pt, densely dashed},
  cellprim/.style={fill=cellP!16, draw=cellP!75!black, thick},
  cellconv/.style={fill=cellC!10, draw=cellC, thick},
  cellbad/.style={fill=bad!14, draw=bad, thick},
  cellconvline/.style={draw=cellC, thick, densely dashed},
  cellshift/.style={draw=cellS, thick, dashed},
  tri/.style={fill=vB!14, draw=vB!55, thin},
  tri0/.style={fill=vB!45, draw=vB!80!black, thick},
}

% ---------------------------------------------------------- titoli e indice
\usepackage{titlesec}
\titleformat{\section}{\spazio{9\baselineskip}\Large\bfseries\color{csez}}{\thesection}{0.8em}{}%
  [\vspace{-0.6ex}{\color{csez}\titlerule[0.7pt]}]
\titleformat{\subsection}{\spazio{8\baselineskip}\large\bfseries\color{csub}}{\thesubsection}{0.8em}{}
\titleformat{\subsubsection}{\spazio{4\baselineskip}\normalsize\bfseries\color{csub}}{}{0pt}{}
\titlespacing*{\section}{0pt}{22pt plus 4pt minus 2pt}{10pt plus 2pt}
\titlespacing*{\subsection}{0pt}{16pt plus 3pt minus 2pt}{6pt plus 2pt}
\titlespacing*{\subsubsection}{0pt}{10pt plus 2pt}{3pt}
\setcounter{tocdepth}{2}
\usepackage[titles]{tocloft}          % titolo «Indice» con lo stile di \section*
\renewcommand{\cftsecfont}{\bfseries}
\renewcommand{\cftsecpagefont}{\bfseries}
\setlength{\cftbeforesecskip}{11pt}
\setlength{\cftbeforesubsecskip}{3pt}

\pagestyle{plain}
\usepackage[colorlinks, linkcolor=csub, citecolor=csub, urlcolor=csub]{hyperref}
\hypersetup{pdftitle={\corso{} - \argomento}, pdfauthor={}}

% ---------------------------------------------------------- etichette numerate
%  Definizione, Risultato, Osservazione (e i riquadri Teorema, Modello) condividono il
%  contatore: Definizione 1.2, Osservazione 1.3, ...  (numerazione per sezione)
%  NB: \if...\else...\fi solo DENTRO le macro (\tit, \sezvoce): un \else scritto
%  dentro \textbf{...} viene accoppiato con l'\ifmmode di \textbf («Extra \else»).
\newcounter{enun}[section]
\renewcommand{\theenun}{\thesection.\arabic{enun}}
\newcommand{\tit}[1]{\if\relax\detokenize{#1}\relax\else\ (#1)\fi}   % « (titolo)» solo se c'è
\newcommand{\etinum}[3]{% colore, nome, titolo facoltativo
  \refstepcounter{enun}\textbf{\color{#1}#2 \theenun\tit{#3}.}\ \ignorespaces}
\newcommand{\etisemp}[3]{% colore, nome, titolo facoltativo (senza numero)
  \textbf{\color{#1}#2\tit{#3}.}\ \ignorespaces}
\newenvironment{definizione}[1][]{\par\addvspace{4pt}\spazio{3\baselineskip}\noindent\etinum{cdef}{Definizione}{#1}}{\par\addvspace{4pt}}
\newenvironment{risultato}[1][]{\par\addvspace{4pt}\spazio{3\baselineskip}\noindent\etinum{cris}{Risultato}{#1}}{\par\addvspace{4pt}}
\newenvironment{osservazione}[1][]{\par\addvspace{4pt}\spazio{3\baselineskip}\noindent\etinum{coss}{Osservazione}{#1}}{\par\addvspace{4pt}}
\newenvironment{interpretazione}[1][]{\par\addvspace{4pt}\spazio{3\baselineskip}\noindent\etisemp{cint}{Interpretazione fisica}{#1}}{\par\addvspace{4pt}}
\newenvironment{riepilogo}[1][]{\par\addvspace{4pt}\spazio{3\baselineskip}\noindent\etisemp{crie}{Riepilogo}{#1}}{\par\addvspace{4pt}}
% introduzione: titolo e testo in nero, spaziatura stretta
\newenvironment{introduzione}{\begingroup\colorlet{csez}{black}\colorlet{crie}{black}%
  \section*{Introduzione}\addcontentsline{toc}{section}{Introduzione}%
  \setlength{\parskip}{3pt plus 1pt}\setlist{itemsep=1pt, topsep=2pt, parsep=0pt}%
  \setlength{\abovedisplayskip}{4pt plus 2pt}\setlength{\belowdisplayskip}{4pt plus 2pt}%
  \setlength{\abovedisplayshortskip}{2pt}\setlength{\belowdisplayshortskip}{3pt}}%
  {\par\endgroup}
% compatibilità con le dispense precedenti
\newenvironment{attenzione}[1][]{\begin{osservazione}[#1]}{\end{osservazione}}

% ---------------------------------------------------------- dimostrazioni a passi
\usepackage[skins, breakable]{tcolorbox}
\newcounter{passo}
\newcommand{\passo}[1]{\par\addvspace{5pt}\stepcounter{passo}%
  \noindent{\bfseries\color{cpasso}Passo \thepasso.}\ \textit{#1}\par\nopagebreak[4]}
\newcommand{\qed}{\ifmmode\qquad\blacksquare\else\unskip\nobreak\ \ensuremath{\blacksquare}\fi}
\newcommand{\iniziobarra}{\setlength{\parskip}{4pt plus 1pt}\setcounter{passo}{0}}
\tcbset{barra/.style={blanker, breakable, parbox=false, left=10pt, right=0pt, top=1pt, bottom=1pt,
  before skip=9pt plus 2pt, after skip=9pt plus 2pt}}
\newtcolorbox{dimostrazione}[1][]{barra, borderline west={1.8pt}{0pt}{cdim},
  before upper={\iniziobarra\textbf{\color{cdim}Dimostrazione\tit{#1}.}\par}}
\newtcolorbox{approfondimento}[1][]{barra, borderline west={1.8pt}{0pt}{cappr},
  before upper={\iniziobarra\textbf{\color{cappr}Approfondimento\tit{#1}.}\ \ignorespaces}}
\newtcolorbox{esempio}[1][]{barra, borderline west={1.8pt}{0pt}{cesm},
  before upper={\iniziobarra\textbf{\color{cesm}Esempio\tit{#1}.}\ \ignorespaces}}
% niente titolo di dimostrazione/approfondimento isolato in fondo alla pagina
\AddToHook{env/dimostrazione/before}{\par\spazio{6\baselineskip}}
\AddToHook{env/approfondimento/before}{\par\spazio{5\baselineskip}}
\AddToHook{env/esempio/before}{\par\spazio{4\baselineskip}}
% conseguenze introdotte dalla domanda a cui rispondono
\newcommand{\domanda}[1]{\par\addvspace{5pt}\spazio{5\baselineskip}\noindent{\color{cris}$\blacktriangleright$}\ \textbf{#1}\par\nopagebreak[4]}
\newcommand{\nsp}[1]{\spazio{#1\baselineskip}}

% ---------------------------------------------------------- riquadri
%  enunciato: riquadro blu con l'enunciato sopra e, dopo \tcblower, il significato
%  \begin{teorema}[di Bloch] enunciato \tcblower significato \end{teorema}
%  (anche: modello). Numerazione comune con Definizione, Osservazione, ...
\newtcolorbox{enunciato}[1]{enhanced, unbreakable, parbox=false,
  colback=cenb, colframe=csez, boxrule=0.6pt, arc=2pt,
  colbacktitle=csez, coltitle=white, fonttitle=\bfseries, title={#1},
  left=9pt, right=9pt, top=6pt, bottom=6pt, toptitle=2pt, bottomtitle=2pt,
  middle=5pt, segmentation style={csez!45, dashed, line width=0.5pt},
  before skip=12pt plus 2pt, after skip=12pt plus 2pt,
  before upper={\setlength{\parskip}{5pt}},
  before lower={\setlength{\parskip}{5pt}\textbf{\color{cint}Cosa dice.}\ \ignorespaces}}
\newenvironment{teorema}[1][]{\par\refstepcounter{enun}\begin{enunciato}{Teorema \theenun\tit{#1}}}{\end{enunciato}}
\newenvironment{modello}[1][]{\par\refstepcounter{enun}\begin{enunciato}{Modello \theenun\tit{#1}}}{\end{enunciato}}
%  chiave: SOLO i risultati essenziali (barra del titolo verde, non spezzabile)
\newtcolorbox{chiave}[1][Risultato]{enhanced, unbreakable, parbox=false,
  colback=cbox, colframe=cris, boxrule=0.6pt, arc=2pt,
  colbacktitle=cris, coltitle=white, fonttitle=\bfseries, title={#1},
  left=9pt, right=9pt, top=5pt, bottom=5pt, toptitle=2pt, bottomtitle=2pt,
  before skip=11pt plus 2pt, after skip=11pt plus 2pt,
  before upper={\setlength{\parskip}{5pt}}}
%  formula da evidenziare senza titolo
\newcommand{\formula}[1]{\par\addvspace{4pt}\begin{center}\begin{tcolorbox}[enhanced, hbox, colback=cform,
  colframe=cformb, boxrule=0.5pt, arc=2pt, left=10pt, right=10pt, top=2pt, bottom=2pt]
  $\displaystyle #1$\end{tcolorbox}\end{center}\par\addvspace{2pt}}
%  informazioni sul corso o curiosità sull'hardware (rarissima)
\newcommand{\titinfo}[1]{\if\relax\detokenize{#1}\relax\else\textbf{#1.}\ \fi}
\newtcolorbox{infobox}[1][]{enhanced, breakable, parbox=false, colback=black!4, frame hidden,
  boxrule=0pt, arc=2pt, left=9pt, right=9pt, top=5pt, bottom=5pt,
  before upper={\setlength{\parskip}{4pt}\titinfo{#1}\ignorespaces}}
%  formulario: riquadro verde spezzabile, una \voce per riga (si spezza tra le voci)
\newtcolorbox{formulario}[1][Formule della lezione]{enhanced, breakable, parbox=false,
  colback=cbox, colframe=cris, boxrule=0.6pt, arc=2pt,
  colbacktitle=cris, coltitle=white, fonttitle=\bfseries, title={#1},
  left=9pt, right=9pt, top=5pt, bottom=7pt, before upper={\setlength{\parskip}{0pt}}}
\newcommand{\gruppo}[1]{\par\addvspace{6pt}\noindent\textbf{#1}\par\nopagebreak\addvspace{4pt}}
\newcommand{\sezvoce}[1]{\if\relax\detokenize{#1}\relax\else\ {\footnotesize\color{black!55}(\S#1)}\fi}
\newcommand{\voce}[3][]{\par\noindent\begin{tabularx}{\linewidth}{@{}>{\raggedright\arraybackslash}p{3.9cm}@{\hspace{8pt}}X@{}}
  #2\sezvoce{#1} & $\displaystyle #3$\end{tabularx}\par\addvspace{4pt}}
%  domande di verifica: \begin{domande}\item \textbf{Domanda?} traccia...\end{domande}
\newlist{domande}{enumerate}{1}
\setlist[domande]{label=\textbf{Q\arabic*.}, leftmargin=2.3em, itemsep=6pt}

% ---------------------------------------------------------- macro generiche
\newcommand{\ingl}[1]{\textit{\color{black!55}(#1)}}   % termine inglese delle slide
\newcommand{\va}{\mathbf a}
\newcommand{\vb}{\mathbf b}
\newcommand{\vk}{\mathbf k}
\newcommand{\vr}{\mathbf r}
\newcommand{\vR}{\mathbf R}
\newcommand{\vG}{\mathbf G}
\newcommand{\vc}{\mathbf c}
\newcommand{\vT}{\mathbf T}
\newcommand{\xh}{\hat{\mathbf x}}
\newcommand{\yh}{\hat{\mathbf y}}
\newcommand{\zh}{\hat{\mathbf z}}
\newcommand{\coord}[1]{\textsf{#1}}
\newcommand{\h}{\tfrac12}
\newcommand{\jj}{\mathrm j}
\newcommand{\ee}{\mathrm e}
\newcommand{\dd}{\mathrm d}
\newcommand{\pd}[2]{\frac{\partial #1}{\partial #2}}
\newcommand{\figura}[1]{\input{figs/#1.tex}}

% ---------------------------------------------------------- float
\renewcommand{\topfraction}{0.9}
\renewcommand{\bottomfraction}{0.8}
\renewcommand{\textfraction}{0.1}
\renewcommand{\floatpagefraction}{0.8}
\setcounter{topnumber}{3}
\setcounter{bottomnumber}{2}
\setcounter{totalnumber}{4}

% ---------------------------------------------------------- aggiunte della Parte 3
\newcommand{\EF}{E_{\mathrm F}}
\newcommand{\kF}{k_{\mathrm F}}
\newcommand{\kB}{k_{\mathrm B}}
\newcommand{\Eg}{E_{\mathrm g}}
\newcommand{\vion}{v_{\mathrm{ion}}}
\newcommand{\Kt}{\tilde K}
\definecolor{occ}{HTML}{2E6FB7}
\definecolor{hole}{HTML}{C0392B}
\definecolor{exc}{HTML}{23884A}
\definecolor{bandfull}{HTML}{8C2A1F}
\definecolor{bandempty}{HTML}{F2E6A0}
\definecolor{bandpart}{HTML}{D9C84A}
%  le tre approssimazioni del §6: stesso riquadro blu di teoremi e modelli
\newenvironment{approssimazione}[1][]{\par\refstepcounter{enun}\begin{enunciato}{Approssimazione \theenun\tit{#1}}}{\end{enunciato}}
\makeatletter
\setlength{\@fptop}{0pt}
\makeatother
 definisci \corso, \argomento, \lezioni.
%  Compilare con pdflatex (scripts/setup_pdflatex_it.sh per la sillabazione).
% ------------------------------------------------------------------
\providecommand{\corso}{Nome del corso}
\providecommand{\argomento}{Argomento}
\providecommand{\lezioni}{Lezioni}
\usepackage[T1]{fontenc}
\usepackage{babel}
\babelprovide[import, main]{italian}
\usepackage{lmodern}            % Latin Modern (se manca: apt-get install -y lmodern)
\usepackage{microtype}
\usepackage[a4paper, left=2.4cm, right=2.4cm, top=2.5cm, bottom=2.7cm]{geometry}
\usepackage[skip=6pt plus 2pt, indent=0pt]{parskip}

\usepackage{amsmath, amssymb}
\usepackage{bm}
\numberwithin{equation}{section}
\usepackage{siunitx}
\sisetup{output-decimal-marker={,}, group-separator={\,}, exponent-product=\cdot}
\DeclareSIUnit{\angstrom}{\text{\AA}}
\usepackage{booktabs, array, tabularx, multirow, longtable}
\usepackage{enumitem}
\setlist{itemsep=3pt, topsep=4pt, parsep=0pt}
\usepackage[font=small, labelfont=bf, format=hang, skip=6pt]{caption}
\usepackage{float}
\usepackage{needspace}
% \spazio{n\baselineskip}: va a capo pagina solo se restano meno di n righe.
% Sostituisce \needspace, il cui \penalty-100 con \vskip elastico fa spezzare
% le pagine prima dei titoli anche quando c'è posto (pagine mezze vuote).
\makeatletter
\newcommand{\spazio}[1]{\par\begingroup\setlength{\dimen@}{#1}%
  \dimen@ii\pagegoal\advance\dimen@ii-\pagetotal
  \ifdim\dimen@>\dimen@ii\ifdim\dimen@ii>\z@\vfil\fi\penalty-10000\fi\endgroup}
\makeatother

% ---------------------------------------------------------- colori
\usepackage{xcolor}
\definecolor{csez}{HTML}{193C82}   % titoli di sezione
\definecolor{csub}{HTML}{3C5995}   % sottosezioni, link
\definecolor{cdef}{HTML}{1E5AAA}   % Definizione
\definecolor{cris}{HTML}{14823C}   % Risultato, riquadro verde
\definecolor{cbox}{HTML}{F1F7F3}   % fondo del riquadro verde
\definecolor{cform}{HTML}{E4F0E8}  % fondo di \formula
\definecolor{cformb}{HTML}{ACD3BA} % bordo di \formula
\definecolor{coss}{HTML}{BE1E1E}   % Osservazione (refusi, convenzioni)
\definecolor{cint}{HTML}{007D87}   % Interpretazione fisica, «Cosa dice»
\definecolor{cappr}{HTML}{783296}  % Approfondimento
\definecolor{crie}{HTML}{C86E00}   % Riepilogo
\definecolor{cdim}{HTML}{6E6E6E}   % Dimostrazione
\definecolor{cpasso}{HTML}{585858} % Passo n.
\definecolor{cesm}{HTML}{4A6FA5}   % Esempio
\definecolor{cenb}{HTML}{EEF3FA}   % fondo del riquadro blu dell'enunciato
% colori per le figure
\definecolor{accent}{HTML}{193C82}
\definecolor{vA}{HTML}{C0392B}
\definecolor{vB}{HTML}{2E6FB7}
\definecolor{vC}{HTML}{23884A}
\definecolor{atomA}{HTML}{E15759}
\definecolor{atomB}{HTML}{4E79A7}
\definecolor{atomC}{HTML}{F28E2B}
\definecolor{atomD}{HTML}{59A14F}
\definecolor{atomE}{HTML}{B07AA1}
\definecolor{cellP}{HTML}{3C9D4E}
\definecolor{cellC}{HTML}{1F4E8C}
\definecolor{cellS}{HTML}{7B3FA0}
\definecolor{bad}{HTML}{C0392B}
\definecolor{Na}{HTML}{F28E2B}
\definecolor{Cl}{HTML}{4E79A7}
% nomi del preambolo precedente (figure e snippet esistenti)
\colorlet{boxdef}{cdef}\colorlet{boxkey}{cris}\colorlet{boxwarn}{coss}
\colorlet{boxex}{black!65}\colorlet{boxdeep}{cappr}\colorlet{mapgray}{black!60}

% ---------------------------------------------------------- figure
\usepackage{tikz}
\usetikzlibrary{arrows.meta, calc, decorations.pathreplacing, patterns}
\usepackage{pgfplots}
\pgfplotsset{compat=1.18, every axis/.append style={font=\footnotesize}, trig format plots=rad}
\usepackage[europeanresistors, cuteinductors, europeanvoltages, straightvoltages]{circuitikz}
\ctikzset{bipoles/length=1.0cm}
\usepackage{tikz-3dplot}
\tikzset{
  vec/.style={-{Stealth[length=2.4mm, width=1.7mm]}, thick},
  bond/.style={black!22, line width=0.5pt},
  latline/.style={black!30, thin},
  edge/.style={black!75, line width=0.7pt},
  edgehid/.style={black!55, line width=0.6pt, densely dashed},
  edgelight/.style={black!50, line width=0.6pt},
  edgelighthid/.style={black!40, line width=0.5pt, densely dashed},
  cellprim/.style={fill=cellP!16, draw=cellP!75!black, thick},
  cellconv/.style={fill=cellC!10, draw=cellC, thick},
  cellbad/.style={fill=bad!14, draw=bad, thick},
  cellconvline/.style={draw=cellC, thick, densely dashed},
  cellshift/.style={draw=cellS, thick, dashed},
  tri/.style={fill=vB!14, draw=vB!55, thin},
  tri0/.style={fill=vB!45, draw=vB!80!black, thick},
}

% ---------------------------------------------------------- titoli e indice
\usepackage{titlesec}
\titleformat{\section}{\spazio{9\baselineskip}\Large\bfseries\color{csez}}{\thesection}{0.8em}{}%
  [\vspace{-0.6ex}{\color{csez}\titlerule[0.7pt]}]
\titleformat{\subsection}{\spazio{8\baselineskip}\large\bfseries\color{csub}}{\thesubsection}{0.8em}{}
\titleformat{\subsubsection}{\spazio{4\baselineskip}\normalsize\bfseries\color{csub}}{}{0pt}{}
\titlespacing*{\section}{0pt}{22pt plus 4pt minus 2pt}{10pt plus 2pt}
\titlespacing*{\subsection}{0pt}{16pt plus 3pt minus 2pt}{6pt plus 2pt}
\titlespacing*{\subsubsection}{0pt}{10pt plus 2pt}{3pt}
\setcounter{tocdepth}{2}
\usepackage[titles]{tocloft}          % titolo «Indice» con lo stile di \section*
\renewcommand{\cftsecfont}{\bfseries}
\renewcommand{\cftsecpagefont}{\bfseries}
\setlength{\cftbeforesecskip}{11pt}
\setlength{\cftbeforesubsecskip}{3pt}

\pagestyle{plain}
\usepackage[colorlinks, linkcolor=csub, citecolor=csub, urlcolor=csub]{hyperref}
\hypersetup{pdftitle={\corso{} - \argomento}, pdfauthor={}}

% ---------------------------------------------------------- etichette numerate
%  Definizione, Risultato, Osservazione (e i riquadri Teorema, Modello) condividono il
%  contatore: Definizione 1.2, Osservazione 1.3, ...  (numerazione per sezione)
%  NB: \if...\else...\fi solo DENTRO le macro (\tit, \sezvoce): un \else scritto
%  dentro \textbf{...} viene accoppiato con l'\ifmmode di \textbf («Extra \else»).
\newcounter{enun}[section]
\renewcommand{\theenun}{\thesection.\arabic{enun}}
\newcommand{\tit}[1]{\if\relax\detokenize{#1}\relax\else\ (#1)\fi}   % « (titolo)» solo se c'è
\newcommand{\etinum}[3]{% colore, nome, titolo facoltativo
  \refstepcounter{enun}\textbf{\color{#1}#2 \theenun\tit{#3}.}\ \ignorespaces}
\newcommand{\etisemp}[3]{% colore, nome, titolo facoltativo (senza numero)
  \textbf{\color{#1}#2\tit{#3}.}\ \ignorespaces}
\newenvironment{definizione}[1][]{\par\addvspace{4pt}\spazio{3\baselineskip}\noindent\etinum{cdef}{Definizione}{#1}}{\par\addvspace{4pt}}
\newenvironment{risultato}[1][]{\par\addvspace{4pt}\spazio{3\baselineskip}\noindent\etinum{cris}{Risultato}{#1}}{\par\addvspace{4pt}}
\newenvironment{osservazione}[1][]{\par\addvspace{4pt}\spazio{3\baselineskip}\noindent\etinum{coss}{Osservazione}{#1}}{\par\addvspace{4pt}}
\newenvironment{interpretazione}[1][]{\par\addvspace{4pt}\spazio{3\baselineskip}\noindent\etisemp{cint}{Interpretazione fisica}{#1}}{\par\addvspace{4pt}}
\newenvironment{riepilogo}[1][]{\par\addvspace{4pt}\spazio{3\baselineskip}\noindent\etisemp{crie}{Riepilogo}{#1}}{\par\addvspace{4pt}}
% introduzione: titolo e testo in nero, spaziatura stretta
\newenvironment{introduzione}{\begingroup\colorlet{csez}{black}\colorlet{crie}{black}%
  \section*{Introduzione}\addcontentsline{toc}{section}{Introduzione}%
  \setlength{\parskip}{3pt plus 1pt}\setlist{itemsep=1pt, topsep=2pt, parsep=0pt}%
  \setlength{\abovedisplayskip}{4pt plus 2pt}\setlength{\belowdisplayskip}{4pt plus 2pt}%
  \setlength{\abovedisplayshortskip}{2pt}\setlength{\belowdisplayshortskip}{3pt}}%
  {\par\endgroup}
% compatibilità con le dispense precedenti
\newenvironment{attenzione}[1][]{\begin{osservazione}[#1]}{\end{osservazione}}

% ---------------------------------------------------------- dimostrazioni a passi
\usepackage[skins, breakable]{tcolorbox}
\newcounter{passo}
\newcommand{\passo}[1]{\par\addvspace{5pt}\stepcounter{passo}%
  \noindent{\bfseries\color{cpasso}Passo \thepasso.}\ \textit{#1}\par\nopagebreak[4]}
\newcommand{\qed}{\ifmmode\qquad\blacksquare\else\unskip\nobreak\ \ensuremath{\blacksquare}\fi}
\newcommand{\iniziobarra}{\setlength{\parskip}{4pt plus 1pt}\setcounter{passo}{0}}
\tcbset{barra/.style={blanker, breakable, parbox=false, left=10pt, right=0pt, top=1pt, bottom=1pt,
  before skip=9pt plus 2pt, after skip=9pt plus 2pt}}
\newtcolorbox{dimostrazione}[1][]{barra, borderline west={1.8pt}{0pt}{cdim},
  before upper={\iniziobarra\textbf{\color{cdim}Dimostrazione\tit{#1}.}\par}}
\newtcolorbox{approfondimento}[1][]{barra, borderline west={1.8pt}{0pt}{cappr},
  before upper={\iniziobarra\textbf{\color{cappr}Approfondimento\tit{#1}.}\ \ignorespaces}}
\newtcolorbox{esempio}[1][]{barra, borderline west={1.8pt}{0pt}{cesm},
  before upper={\iniziobarra\textbf{\color{cesm}Esempio\tit{#1}.}\ \ignorespaces}}
% niente titolo di dimostrazione/approfondimento isolato in fondo alla pagina
\AddToHook{env/dimostrazione/before}{\par\spazio{6\baselineskip}}
\AddToHook{env/approfondimento/before}{\par\spazio{5\baselineskip}}
\AddToHook{env/esempio/before}{\par\spazio{4\baselineskip}}
% conseguenze introdotte dalla domanda a cui rispondono
\newcommand{\domanda}[1]{\par\addvspace{5pt}\spazio{5\baselineskip}\noindent{\color{cris}$\blacktriangleright$}\ \textbf{#1}\par\nopagebreak[4]}
\newcommand{\nsp}[1]{\spazio{#1\baselineskip}}

% ---------------------------------------------------------- riquadri
%  enunciato: riquadro blu con l'enunciato sopra e, dopo \tcblower, il significato
%  \begin{teorema}[di Bloch] enunciato \tcblower significato \end{teorema}
%  (anche: modello). Numerazione comune con Definizione, Osservazione, ...
\newtcolorbox{enunciato}[1]{enhanced, unbreakable, parbox=false,
  colback=cenb, colframe=csez, boxrule=0.6pt, arc=2pt,
  colbacktitle=csez, coltitle=white, fonttitle=\bfseries, title={#1},
  left=9pt, right=9pt, top=6pt, bottom=6pt, toptitle=2pt, bottomtitle=2pt,
  middle=5pt, segmentation style={csez!45, dashed, line width=0.5pt},
  before skip=12pt plus 2pt, after skip=12pt plus 2pt,
  before upper={\setlength{\parskip}{5pt}},
  before lower={\setlength{\parskip}{5pt}\textbf{\color{cint}Cosa dice.}\ \ignorespaces}}
\newenvironment{teorema}[1][]{\par\refstepcounter{enun}\begin{enunciato}{Teorema \theenun\tit{#1}}}{\end{enunciato}}
\newenvironment{modello}[1][]{\par\refstepcounter{enun}\begin{enunciato}{Modello \theenun\tit{#1}}}{\end{enunciato}}
%  chiave: SOLO i risultati essenziali (barra del titolo verde, non spezzabile)
\newtcolorbox{chiave}[1][Risultato]{enhanced, unbreakable, parbox=false,
  colback=cbox, colframe=cris, boxrule=0.6pt, arc=2pt,
  colbacktitle=cris, coltitle=white, fonttitle=\bfseries, title={#1},
  left=9pt, right=9pt, top=5pt, bottom=5pt, toptitle=2pt, bottomtitle=2pt,
  before skip=11pt plus 2pt, after skip=11pt plus 2pt,
  before upper={\setlength{\parskip}{5pt}}}
%  formula da evidenziare senza titolo
\newcommand{\formula}[1]{\par\addvspace{4pt}\begin{center}\begin{tcolorbox}[enhanced, hbox, colback=cform,
  colframe=cformb, boxrule=0.5pt, arc=2pt, left=10pt, right=10pt, top=2pt, bottom=2pt]
  $\displaystyle #1$\end{tcolorbox}\end{center}\par\addvspace{2pt}}
%  informazioni sul corso o curiosità sull'hardware (rarissima)
\newcommand{\titinfo}[1]{\if\relax\detokenize{#1}\relax\else\textbf{#1.}\ \fi}
\newtcolorbox{infobox}[1][]{enhanced, breakable, parbox=false, colback=black!4, frame hidden,
  boxrule=0pt, arc=2pt, left=9pt, right=9pt, top=5pt, bottom=5pt,
  before upper={\setlength{\parskip}{4pt}\titinfo{#1}\ignorespaces}}
%  formulario: riquadro verde spezzabile, una \voce per riga (si spezza tra le voci)
\newtcolorbox{formulario}[1][Formule della lezione]{enhanced, breakable, parbox=false,
  colback=cbox, colframe=cris, boxrule=0.6pt, arc=2pt,
  colbacktitle=cris, coltitle=white, fonttitle=\bfseries, title={#1},
  left=9pt, right=9pt, top=5pt, bottom=7pt, before upper={\setlength{\parskip}{0pt}}}
\newcommand{\gruppo}[1]{\par\addvspace{6pt}\noindent\textbf{#1}\par\nopagebreak\addvspace{4pt}}
\newcommand{\sezvoce}[1]{\if\relax\detokenize{#1}\relax\else\ {\footnotesize\color{black!55}(\S#1)}\fi}
\newcommand{\voce}[3][]{\par\noindent\begin{tabularx}{\linewidth}{@{}>{\raggedright\arraybackslash}p{3.9cm}@{\hspace{8pt}}X@{}}
  #2\sezvoce{#1} & $\displaystyle #3$\end{tabularx}\par\addvspace{4pt}}
%  domande di verifica: \begin{domande}\item \textbf{Domanda?} traccia...\end{domande}
\newlist{domande}{enumerate}{1}
\setlist[domande]{label=\textbf{Q\arabic*.}, leftmargin=2.3em, itemsep=6pt}

% ---------------------------------------------------------- macro generiche
\newcommand{\ingl}[1]{\textit{\color{black!55}(#1)}}   % termine inglese delle slide
\newcommand{\va}{\mathbf a}
\newcommand{\vb}{\mathbf b}
\newcommand{\vk}{\mathbf k}
\newcommand{\vr}{\mathbf r}
\newcommand{\vR}{\mathbf R}
\newcommand{\vG}{\mathbf G}
\newcommand{\vc}{\mathbf c}
\newcommand{\vT}{\mathbf T}
\newcommand{\xh}{\hat{\mathbf x}}
\newcommand{\yh}{\hat{\mathbf y}}
\newcommand{\zh}{\hat{\mathbf z}}
\newcommand{\coord}[1]{\textsf{#1}}
\newcommand{\h}{\tfrac12}
\newcommand{\jj}{\mathrm j}
\newcommand{\ee}{\mathrm e}
\newcommand{\dd}{\mathrm d}
\newcommand{\pd}[2]{\frac{\partial #1}{\partial #2}}
\newcommand{\figura}[1]{\input{figs/#1.tex}}

% ---------------------------------------------------------- float
\renewcommand{\topfraction}{0.9}
\renewcommand{\bottomfraction}{0.8}
\renewcommand{\textfraction}{0.1}
\renewcommand{\floatpagefraction}{0.8}
\setcounter{topnumber}{3}
\setcounter{bottomnumber}{2}
\setcounter{totalnumber}{4}

% ---------------------------------------------------------- aggiunte della Parte 3
\newcommand{\EF}{E_{\mathrm F}}
\newcommand{\kF}{k_{\mathrm F}}
\newcommand{\kB}{k_{\mathrm B}}
\newcommand{\Eg}{E_{\mathrm g}}
\newcommand{\vion}{v_{\mathrm{ion}}}
\newcommand{\Kt}{\tilde K}
\definecolor{occ}{HTML}{2E6FB7}
\definecolor{hole}{HTML}{C0392B}
\definecolor{exc}{HTML}{23884A}
\definecolor{bandfull}{HTML}{8C2A1F}
\definecolor{bandempty}{HTML}{F2E6A0}
\definecolor{bandpart}{HTML}{D9C84A}
%  le tre approssimazioni del §6: stesso riquadro blu di teoremi e modelli
\newenvironment{approssimazione}[1][]{\par\refstepcounter{enun}\begin{enunciato}{Approssimazione \theenun\tit{#1}}}{\end{enunciato}}
\makeatletter
\setlength{\@fptop}{0pt}
\makeatother
 definisci \corso, \argomento, \lezioni.
%  Compilare con pdflatex (scripts/setup_pdflatex_it.sh per la sillabazione).
% ------------------------------------------------------------------
\providecommand{\corso}{Nome del corso}
\providecommand{\argomento}{Argomento}
\providecommand{\lezioni}{Lezioni}
\usepackage[T1]{fontenc}
\usepackage{babel}
\babelprovide[import, main]{italian}
\usepackage{lmodern}            % Latin Modern (se manca: apt-get install -y lmodern)
\usepackage{microtype}
\usepackage[a4paper, left=2.4cm, right=2.4cm, top=2.5cm, bottom=2.7cm]{geometry}
\usepackage[skip=6pt plus 2pt, indent=0pt]{parskip}

\usepackage{amsmath, amssymb}
\usepackage{bm}
\numberwithin{equation}{section}
\usepackage{siunitx}
\sisetup{output-decimal-marker={,}, group-separator={\,}, exponent-product=\cdot}
\DeclareSIUnit{\angstrom}{\text{\AA}}
\usepackage{booktabs, array, tabularx, multirow, longtable}
\usepackage{enumitem}
\setlist{itemsep=3pt, topsep=4pt, parsep=0pt}
\usepackage[font=small, labelfont=bf, format=hang, skip=6pt]{caption}
\usepackage{float}
\usepackage{needspace}
% \spazio{n\baselineskip}: va a capo pagina solo se restano meno di n righe.
% Sostituisce \needspace, il cui \penalty-100 con \vskip elastico fa spezzare
% le pagine prima dei titoli anche quando c'è posto (pagine mezze vuote).
\makeatletter
\newcommand{\spazio}[1]{\par\begingroup\setlength{\dimen@}{#1}%
  \dimen@ii\pagegoal\advance\dimen@ii-\pagetotal
  \ifdim\dimen@>\dimen@ii\ifdim\dimen@ii>\z@\vfil\fi\penalty-10000\fi\endgroup}
\makeatother

% ---------------------------------------------------------- colori
\usepackage{xcolor}
\definecolor{csez}{HTML}{193C82}   % titoli di sezione
\definecolor{csub}{HTML}{3C5995}   % sottosezioni, link
\definecolor{cdef}{HTML}{1E5AAA}   % Definizione
\definecolor{cris}{HTML}{14823C}   % Risultato, riquadro verde
\definecolor{cbox}{HTML}{F1F7F3}   % fondo del riquadro verde
\definecolor{cform}{HTML}{E4F0E8}  % fondo di \formula
\definecolor{cformb}{HTML}{ACD3BA} % bordo di \formula
\definecolor{coss}{HTML}{BE1E1E}   % Osservazione (refusi, convenzioni)
\definecolor{cint}{HTML}{007D87}   % Interpretazione fisica, «Cosa dice»
\definecolor{cappr}{HTML}{783296}  % Approfondimento
\definecolor{crie}{HTML}{C86E00}   % Riepilogo
\definecolor{cdim}{HTML}{6E6E6E}   % Dimostrazione
\definecolor{cpasso}{HTML}{585858} % Passo n.
\definecolor{cesm}{HTML}{4A6FA5}   % Esempio
\definecolor{cenb}{HTML}{EEF3FA}   % fondo del riquadro blu dell'enunciato
% colori per le figure
\definecolor{accent}{HTML}{193C82}
\definecolor{vA}{HTML}{C0392B}
\definecolor{vB}{HTML}{2E6FB7}
\definecolor{vC}{HTML}{23884A}
\definecolor{atomA}{HTML}{E15759}
\definecolor{atomB}{HTML}{4E79A7}
\definecolor{atomC}{HTML}{F28E2B}
\definecolor{atomD}{HTML}{59A14F}
\definecolor{atomE}{HTML}{B07AA1}
\definecolor{cellP}{HTML}{3C9D4E}
\definecolor{cellC}{HTML}{1F4E8C}
\definecolor{cellS}{HTML}{7B3FA0}
\definecolor{bad}{HTML}{C0392B}
\definecolor{Na}{HTML}{F28E2B}
\definecolor{Cl}{HTML}{4E79A7}
% nomi del preambolo precedente (figure e snippet esistenti)
\colorlet{boxdef}{cdef}\colorlet{boxkey}{cris}\colorlet{boxwarn}{coss}
\colorlet{boxex}{black!65}\colorlet{boxdeep}{cappr}\colorlet{mapgray}{black!60}

% ---------------------------------------------------------- figure
\usepackage{tikz}
\usetikzlibrary{arrows.meta, calc, decorations.pathreplacing, patterns}
\usepackage{pgfplots}
\pgfplotsset{compat=1.18, every axis/.append style={font=\footnotesize}, trig format plots=rad}
\usepackage[europeanresistors, cuteinductors, europeanvoltages, straightvoltages]{circuitikz}
\ctikzset{bipoles/length=1.0cm}
\usepackage{tikz-3dplot}
\tikzset{
  vec/.style={-{Stealth[length=2.4mm, width=1.7mm]}, thick},
  bond/.style={black!22, line width=0.5pt},
  latline/.style={black!30, thin},
  edge/.style={black!75, line width=0.7pt},
  edgehid/.style={black!55, line width=0.6pt, densely dashed},
  edgelight/.style={black!50, line width=0.6pt},
  edgelighthid/.style={black!40, line width=0.5pt, densely dashed},
  cellprim/.style={fill=cellP!16, draw=cellP!75!black, thick},
  cellconv/.style={fill=cellC!10, draw=cellC, thick},
  cellbad/.style={fill=bad!14, draw=bad, thick},
  cellconvline/.style={draw=cellC, thick, densely dashed},
  cellshift/.style={draw=cellS, thick, dashed},
  tri/.style={fill=vB!14, draw=vB!55, thin},
  tri0/.style={fill=vB!45, draw=vB!80!black, thick},
}

% ---------------------------------------------------------- titoli e indice
\usepackage{titlesec}
\titleformat{\section}{\spazio{9\baselineskip}\Large\bfseries\color{csez}}{\thesection}{0.8em}{}%
  [\vspace{-0.6ex}{\color{csez}\titlerule[0.7pt]}]
\titleformat{\subsection}{\spazio{8\baselineskip}\large\bfseries\color{csub}}{\thesubsection}{0.8em}{}
\titleformat{\subsubsection}{\spazio{4\baselineskip}\normalsize\bfseries\color{csub}}{}{0pt}{}
\titlespacing*{\section}{0pt}{22pt plus 4pt minus 2pt}{10pt plus 2pt}
\titlespacing*{\subsection}{0pt}{16pt plus 3pt minus 2pt}{6pt plus 2pt}
\titlespacing*{\subsubsection}{0pt}{10pt plus 2pt}{3pt}
\setcounter{tocdepth}{2}
\usepackage[titles]{tocloft}          % titolo «Indice» con lo stile di \section*
\renewcommand{\cftsecfont}{\bfseries}
\renewcommand{\cftsecpagefont}{\bfseries}
\setlength{\cftbeforesecskip}{11pt}
\setlength{\cftbeforesubsecskip}{3pt}

\pagestyle{plain}
\usepackage[colorlinks, linkcolor=csub, citecolor=csub, urlcolor=csub]{hyperref}
\hypersetup{pdftitle={\corso{} - \argomento}, pdfauthor={}}

% ---------------------------------------------------------- etichette numerate
%  Definizione, Risultato, Osservazione (e i riquadri Teorema, Modello) condividono il
%  contatore: Definizione 1.2, Osservazione 1.3, ...  (numerazione per sezione)
%  NB: \if...\else...\fi solo DENTRO le macro (\tit, \sezvoce): un \else scritto
%  dentro \textbf{...} viene accoppiato con l'\ifmmode di \textbf («Extra \else»).
\newcounter{enun}[section]
\renewcommand{\theenun}{\thesection.\arabic{enun}}
\newcommand{\tit}[1]{\if\relax\detokenize{#1}\relax\else\ (#1)\fi}   % « (titolo)» solo se c'è
\newcommand{\etinum}[3]{% colore, nome, titolo facoltativo
  \refstepcounter{enun}\textbf{\color{#1}#2 \theenun\tit{#3}.}\ \ignorespaces}
\newcommand{\etisemp}[3]{% colore, nome, titolo facoltativo (senza numero)
  \textbf{\color{#1}#2\tit{#3}.}\ \ignorespaces}
\newenvironment{definizione}[1][]{\par\addvspace{4pt}\spazio{3\baselineskip}\noindent\etinum{cdef}{Definizione}{#1}}{\par\addvspace{4pt}}
\newenvironment{risultato}[1][]{\par\addvspace{4pt}\spazio{3\baselineskip}\noindent\etinum{cris}{Risultato}{#1}}{\par\addvspace{4pt}}
\newenvironment{osservazione}[1][]{\par\addvspace{4pt}\spazio{3\baselineskip}\noindent\etinum{coss}{Osservazione}{#1}}{\par\addvspace{4pt}}
\newenvironment{interpretazione}[1][]{\par\addvspace{4pt}\spazio{3\baselineskip}\noindent\etisemp{cint}{Interpretazione fisica}{#1}}{\par\addvspace{4pt}}
\newenvironment{riepilogo}[1][]{\par\addvspace{4pt}\spazio{3\baselineskip}\noindent\etisemp{crie}{Riepilogo}{#1}}{\par\addvspace{4pt}}
% introduzione: titolo e testo in nero, spaziatura stretta
\newenvironment{introduzione}{\begingroup\colorlet{csez}{black}\colorlet{crie}{black}%
  \section*{Introduzione}\addcontentsline{toc}{section}{Introduzione}%
  \setlength{\parskip}{3pt plus 1pt}\setlist{itemsep=1pt, topsep=2pt, parsep=0pt}%
  \setlength{\abovedisplayskip}{4pt plus 2pt}\setlength{\belowdisplayskip}{4pt plus 2pt}%
  \setlength{\abovedisplayshortskip}{2pt}\setlength{\belowdisplayshortskip}{3pt}}%
  {\par\endgroup}
% compatibilità con le dispense precedenti
\newenvironment{attenzione}[1][]{\begin{osservazione}[#1]}{\end{osservazione}}

% ---------------------------------------------------------- dimostrazioni a passi
\usepackage[skins, breakable]{tcolorbox}
\newcounter{passo}
\newcommand{\passo}[1]{\par\addvspace{5pt}\stepcounter{passo}%
  \noindent{\bfseries\color{cpasso}Passo \thepasso.}\ \textit{#1}\par\nopagebreak[4]}
\newcommand{\qed}{\ifmmode\qquad\blacksquare\else\unskip\nobreak\ \ensuremath{\blacksquare}\fi}
\newcommand{\iniziobarra}{\setlength{\parskip}{4pt plus 1pt}\setcounter{passo}{0}}
\tcbset{barra/.style={blanker, breakable, parbox=false, left=10pt, right=0pt, top=1pt, bottom=1pt,
  before skip=9pt plus 2pt, after skip=9pt plus 2pt}}
\newtcolorbox{dimostrazione}[1][]{barra, borderline west={1.8pt}{0pt}{cdim},
  before upper={\iniziobarra\textbf{\color{cdim}Dimostrazione\tit{#1}.}\par}}
\newtcolorbox{approfondimento}[1][]{barra, borderline west={1.8pt}{0pt}{cappr},
  before upper={\iniziobarra\textbf{\color{cappr}Approfondimento\tit{#1}.}\ \ignorespaces}}
\newtcolorbox{esempio}[1][]{barra, borderline west={1.8pt}{0pt}{cesm},
  before upper={\iniziobarra\textbf{\color{cesm}Esempio\tit{#1}.}\ \ignorespaces}}
% niente titolo di dimostrazione/approfondimento isolato in fondo alla pagina
\AddToHook{env/dimostrazione/before}{\par\spazio{6\baselineskip}}
\AddToHook{env/approfondimento/before}{\par\spazio{5\baselineskip}}
\AddToHook{env/esempio/before}{\par\spazio{4\baselineskip}}
% conseguenze introdotte dalla domanda a cui rispondono
\newcommand{\domanda}[1]{\par\addvspace{5pt}\spazio{5\baselineskip}\noindent{\color{cris}$\blacktriangleright$}\ \textbf{#1}\par\nopagebreak[4]}
\newcommand{\nsp}[1]{\spazio{#1\baselineskip}}

% ---------------------------------------------------------- riquadri
%  enunciato: riquadro blu con l'enunciato sopra e, dopo \tcblower, il significato
%  \begin{teorema}[di Bloch] enunciato \tcblower significato \end{teorema}
%  (anche: modello). Numerazione comune con Definizione, Osservazione, ...
\newtcolorbox{enunciato}[1]{enhanced, unbreakable, parbox=false,
  colback=cenb, colframe=csez, boxrule=0.6pt, arc=2pt,
  colbacktitle=csez, coltitle=white, fonttitle=\bfseries, title={#1},
  left=9pt, right=9pt, top=6pt, bottom=6pt, toptitle=2pt, bottomtitle=2pt,
  middle=5pt, segmentation style={csez!45, dashed, line width=0.5pt},
  before skip=12pt plus 2pt, after skip=12pt plus 2pt,
  before upper={\setlength{\parskip}{5pt}},
  before lower={\setlength{\parskip}{5pt}\textbf{\color{cint}Cosa dice.}\ \ignorespaces}}
\newenvironment{teorema}[1][]{\par\refstepcounter{enun}\begin{enunciato}{Teorema \theenun\tit{#1}}}{\end{enunciato}}
\newenvironment{modello}[1][]{\par\refstepcounter{enun}\begin{enunciato}{Modello \theenun\tit{#1}}}{\end{enunciato}}
%  chiave: SOLO i risultati essenziali (barra del titolo verde, non spezzabile)
\newtcolorbox{chiave}[1][Risultato]{enhanced, unbreakable, parbox=false,
  colback=cbox, colframe=cris, boxrule=0.6pt, arc=2pt,
  colbacktitle=cris, coltitle=white, fonttitle=\bfseries, title={#1},
  left=9pt, right=9pt, top=5pt, bottom=5pt, toptitle=2pt, bottomtitle=2pt,
  before skip=11pt plus 2pt, after skip=11pt plus 2pt,
  before upper={\setlength{\parskip}{5pt}}}
%  formula da evidenziare senza titolo
\newcommand{\formula}[1]{\par\addvspace{4pt}\begin{center}\begin{tcolorbox}[enhanced, hbox, colback=cform,
  colframe=cformb, boxrule=0.5pt, arc=2pt, left=10pt, right=10pt, top=2pt, bottom=2pt]
  $\displaystyle #1$\end{tcolorbox}\end{center}\par\addvspace{2pt}}
%  informazioni sul corso o curiosità sull'hardware (rarissima)
\newcommand{\titinfo}[1]{\if\relax\detokenize{#1}\relax\else\textbf{#1.}\ \fi}
\newtcolorbox{infobox}[1][]{enhanced, breakable, parbox=false, colback=black!4, frame hidden,
  boxrule=0pt, arc=2pt, left=9pt, right=9pt, top=5pt, bottom=5pt,
  before upper={\setlength{\parskip}{4pt}\titinfo{#1}\ignorespaces}}
%  formulario: riquadro verde spezzabile, una \voce per riga (si spezza tra le voci)
\newtcolorbox{formulario}[1][Formule della lezione]{enhanced, breakable, parbox=false,
  colback=cbox, colframe=cris, boxrule=0.6pt, arc=2pt,
  colbacktitle=cris, coltitle=white, fonttitle=\bfseries, title={#1},
  left=9pt, right=9pt, top=5pt, bottom=7pt, before upper={\setlength{\parskip}{0pt}}}
\newcommand{\gruppo}[1]{\par\addvspace{6pt}\noindent\textbf{#1}\par\nopagebreak\addvspace{4pt}}
\newcommand{\sezvoce}[1]{\if\relax\detokenize{#1}\relax\else\ {\footnotesize\color{black!55}(\S#1)}\fi}
\newcommand{\voce}[3][]{\par\noindent\begin{tabularx}{\linewidth}{@{}>{\raggedright\arraybackslash}p{3.9cm}@{\hspace{8pt}}X@{}}
  #2\sezvoce{#1} & $\displaystyle #3$\end{tabularx}\par\addvspace{4pt}}
%  domande di verifica: \begin{domande}\item \textbf{Domanda?} traccia...\end{domande}
\newlist{domande}{enumerate}{1}
\setlist[domande]{label=\textbf{Q\arabic*.}, leftmargin=2.3em, itemsep=6pt}

% ---------------------------------------------------------- macro generiche
\newcommand{\ingl}[1]{\textit{\color{black!55}(#1)}}   % termine inglese delle slide
\newcommand{\va}{\mathbf a}
\newcommand{\vb}{\mathbf b}
\newcommand{\vk}{\mathbf k}
\newcommand{\vr}{\mathbf r}
\newcommand{\vR}{\mathbf R}
\newcommand{\vG}{\mathbf G}
\newcommand{\vc}{\mathbf c}
\newcommand{\vT}{\mathbf T}
\newcommand{\xh}{\hat{\mathbf x}}
\newcommand{\yh}{\hat{\mathbf y}}
\newcommand{\zh}{\hat{\mathbf z}}
\newcommand{\coord}[1]{\textsf{#1}}
\newcommand{\h}{\tfrac12}
\newcommand{\jj}{\mathrm j}
\newcommand{\ee}{\mathrm e}
\newcommand{\dd}{\mathrm d}
\newcommand{\pd}[2]{\frac{\partial #1}{\partial #2}}
\newcommand{\figura}[1]{\input{figs/#1.tex}}

% ---------------------------------------------------------- float
\renewcommand{\topfraction}{0.9}
\renewcommand{\bottomfraction}{0.8}
\renewcommand{\textfraction}{0.1}
\renewcommand{\floatpagefraction}{0.8}
\setcounter{topnumber}{3}
\setcounter{bottomnumber}{2}
\setcounter{totalnumber}{4}

% ---------------------------------------------------------- aggiunte della Parte 3
\newcommand{\EF}{E_{\mathrm F}}
\newcommand{\kF}{k_{\mathrm F}}
\newcommand{\kB}{k_{\mathrm B}}
\newcommand{\Eg}{E_{\mathrm g}}
\newcommand{\vion}{v_{\mathrm{ion}}}
\newcommand{\Kt}{\tilde K}
\definecolor{occ}{HTML}{2E6FB7}
\definecolor{hole}{HTML}{C0392B}
\definecolor{exc}{HTML}{23884A}
\definecolor{bandfull}{HTML}{8C2A1F}
\definecolor{bandempty}{HTML}{F2E6A0}
\definecolor{bandpart}{HTML}{D9C84A}
%  le tre approssimazioni del §6: stesso riquadro blu di teoremi e modelli
\newenvironment{approssimazione}[1][]{\par\refstepcounter{enun}\begin{enunciato}{Approssimazione \theenun\tit{#1}}}{\end{enunciato}}
\makeatletter
\setlength{\@fptop}{0pt}
\makeatother
 definisci \corso, \argomento, \lezioni.
%  Compilare con pdflatex (scripts/setup_pdflatex_it.sh per la sillabazione).
% ------------------------------------------------------------------
\providecommand{\corso}{Nome del corso}
\providecommand{\argomento}{Argomento}
\providecommand{\lezioni}{Lezioni}
\usepackage[T1]{fontenc}
\usepackage{babel}
\babelprovide[import, main]{italian}
\usepackage{lmodern}            % Latin Modern (se manca: apt-get install -y lmodern)
\usepackage{microtype}
\usepackage[a4paper, left=2.4cm, right=2.4cm, top=2.5cm, bottom=2.7cm]{geometry}
\usepackage[skip=6pt plus 2pt, indent=0pt]{parskip}

\usepackage{amsmath, amssymb}
\usepackage{bm}
\numberwithin{equation}{section}
\usepackage{siunitx}
\sisetup{output-decimal-marker={,}, group-separator={\,}, exponent-product=\cdot}
\DeclareSIUnit{\angstrom}{\text{\AA}}
\usepackage{booktabs, array, tabularx, multirow, longtable}
\usepackage{enumitem}
\setlist{itemsep=3pt, topsep=4pt, parsep=0pt}
\usepackage[font=small, labelfont=bf, format=hang, skip=6pt]{caption}
\usepackage{float}
\usepackage{needspace}
% \spazio{n\baselineskip}: va a capo pagina solo se restano meno di n righe.
% Sostituisce \needspace, il cui \penalty-100 con \vskip elastico fa spezzare
% le pagine prima dei titoli anche quando c'è posto (pagine mezze vuote).
\makeatletter
\newcommand{\spazio}[1]{\par\begingroup\setlength{\dimen@}{#1}%
  \dimen@ii\pagegoal\advance\dimen@ii-\pagetotal
  \ifdim\dimen@>\dimen@ii\ifdim\dimen@ii>\z@\vfil\fi\penalty-10000\fi\endgroup}
\makeatother

% ---------------------------------------------------------- colori
\usepackage{xcolor}
\definecolor{csez}{HTML}{193C82}   % titoli di sezione
\definecolor{csub}{HTML}{3C5995}   % sottosezioni, link
\definecolor{cdef}{HTML}{1E5AAA}   % Definizione
\definecolor{cris}{HTML}{14823C}   % Risultato, riquadro verde
\definecolor{cbox}{HTML}{F1F7F3}   % fondo del riquadro verde
\definecolor{cform}{HTML}{E4F0E8}  % fondo di \formula
\definecolor{cformb}{HTML}{ACD3BA} % bordo di \formula
\definecolor{coss}{HTML}{BE1E1E}   % Osservazione (refusi, convenzioni)
\definecolor{cint}{HTML}{007D87}   % Interpretazione fisica, «Cosa dice»
\definecolor{cappr}{HTML}{783296}  % Approfondimento
\definecolor{crie}{HTML}{C86E00}   % Riepilogo
\definecolor{cdim}{HTML}{6E6E6E}   % Dimostrazione
\definecolor{cpasso}{HTML}{585858} % Passo n.
\definecolor{cesm}{HTML}{4A6FA5}   % Esempio
\definecolor{cenb}{HTML}{EEF3FA}   % fondo del riquadro blu dell'enunciato
% colori per le figure
\definecolor{accent}{HTML}{193C82}
\definecolor{vA}{HTML}{C0392B}
\definecolor{vB}{HTML}{2E6FB7}
\definecolor{vC}{HTML}{23884A}
\definecolor{atomA}{HTML}{E15759}
\definecolor{atomB}{HTML}{4E79A7}
\definecolor{atomC}{HTML}{F28E2B}
\definecolor{atomD}{HTML}{59A14F}
\definecolor{atomE}{HTML}{B07AA1}
\definecolor{cellP}{HTML}{3C9D4E}
\definecolor{cellC}{HTML}{1F4E8C}
\definecolor{cellS}{HTML}{7B3FA0}
\definecolor{bad}{HTML}{C0392B}
\definecolor{Na}{HTML}{F28E2B}
\definecolor{Cl}{HTML}{4E79A7}
% nomi del preambolo precedente (figure e snippet esistenti)
\colorlet{boxdef}{cdef}\colorlet{boxkey}{cris}\colorlet{boxwarn}{coss}
\colorlet{boxex}{black!65}\colorlet{boxdeep}{cappr}\colorlet{mapgray}{black!60}

% ---------------------------------------------------------- figure
\usepackage{tikz}
\usetikzlibrary{arrows.meta, calc, decorations.pathreplacing, patterns}
\usepackage{pgfplots}
\pgfplotsset{compat=1.18, every axis/.append style={font=\footnotesize}, trig format plots=rad}
\usepackage[europeanresistors, cuteinductors, europeanvoltages, straightvoltages]{circuitikz}
\ctikzset{bipoles/length=1.0cm}
\usepackage{tikz-3dplot}
\tikzset{
  vec/.style={-{Stealth[length=2.4mm, width=1.7mm]}, thick},
  bond/.style={black!22, line width=0.5pt},
  latline/.style={black!30, thin},
  edge/.style={black!75, line width=0.7pt},
  edgehid/.style={black!55, line width=0.6pt, densely dashed},
  edgelight/.style={black!50, line width=0.6pt},
  edgelighthid/.style={black!40, line width=0.5pt, densely dashed},
  cellprim/.style={fill=cellP!16, draw=cellP!75!black, thick},
  cellconv/.style={fill=cellC!10, draw=cellC, thick},
  cellbad/.style={fill=bad!14, draw=bad, thick},
  cellconvline/.style={draw=cellC, thick, densely dashed},
  cellshift/.style={draw=cellS, thick, dashed},
  tri/.style={fill=vB!14, draw=vB!55, thin},
  tri0/.style={fill=vB!45, draw=vB!80!black, thick},
}

% ---------------------------------------------------------- titoli e indice
\usepackage{titlesec}
\titleformat{\section}{\spazio{9\baselineskip}\Large\bfseries\color{csez}}{\thesection}{0.8em}{}%
  [\vspace{-0.6ex}{\color{csez}\titlerule[0.7pt]}]
\titleformat{\subsection}{\spazio{8\baselineskip}\large\bfseries\color{csub}}{\thesubsection}{0.8em}{}
\titleformat{\subsubsection}{\spazio{4\baselineskip}\normalsize\bfseries\color{csub}}{}{0pt}{}
\titlespacing*{\section}{0pt}{22pt plus 4pt minus 2pt}{10pt plus 2pt}
\titlespacing*{\subsection}{0pt}{16pt plus 3pt minus 2pt}{6pt plus 2pt}
\titlespacing*{\subsubsection}{0pt}{10pt plus 2pt}{3pt}
\setcounter{tocdepth}{2}
\usepackage[titles]{tocloft}          % titolo «Indice» con lo stile di \section*
\renewcommand{\cftsecfont}{\bfseries}
\renewcommand{\cftsecpagefont}{\bfseries}
\setlength{\cftbeforesecskip}{11pt}
\setlength{\cftbeforesubsecskip}{3pt}

\pagestyle{plain}
\usepackage[colorlinks, linkcolor=csub, citecolor=csub, urlcolor=csub]{hyperref}
\hypersetup{pdftitle={\corso{} - \argomento}, pdfauthor={}}

% ---------------------------------------------------------- etichette numerate
%  Definizione, Risultato, Osservazione (e i riquadri Teorema, Modello) condividono il
%  contatore: Definizione 1.2, Osservazione 1.3, ...  (numerazione per sezione)
%  NB: \if...\else...\fi solo DENTRO le macro (\tit, \sezvoce): un \else scritto
%  dentro \textbf{...} viene accoppiato con l'\ifmmode di \textbf («Extra \else»).
\newcounter{enun}[section]
\renewcommand{\theenun}{\thesection.\arabic{enun}}
\newcommand{\tit}[1]{\if\relax\detokenize{#1}\relax\else\ (#1)\fi}   % « (titolo)» solo se c'è
\newcommand{\etinum}[3]{% colore, nome, titolo facoltativo
  \refstepcounter{enun}\textbf{\color{#1}#2 \theenun\tit{#3}.}\ \ignorespaces}
\newcommand{\etisemp}[3]{% colore, nome, titolo facoltativo (senza numero)
  \textbf{\color{#1}#2\tit{#3}.}\ \ignorespaces}
\newenvironment{definizione}[1][]{\par\addvspace{4pt}\spazio{3\baselineskip}\noindent\etinum{cdef}{Definizione}{#1}}{\par\addvspace{4pt}}
\newenvironment{risultato}[1][]{\par\addvspace{4pt}\spazio{3\baselineskip}\noindent\etinum{cris}{Risultato}{#1}}{\par\addvspace{4pt}}
\newenvironment{osservazione}[1][]{\par\addvspace{4pt}\spazio{3\baselineskip}\noindent\etinum{coss}{Osservazione}{#1}}{\par\addvspace{4pt}}
\newenvironment{interpretazione}[1][]{\par\addvspace{4pt}\spazio{3\baselineskip}\noindent\etisemp{cint}{Interpretazione fisica}{#1}}{\par\addvspace{4pt}}
\newenvironment{riepilogo}[1][]{\par\addvspace{4pt}\spazio{3\baselineskip}\noindent\etisemp{crie}{Riepilogo}{#1}}{\par\addvspace{4pt}}
% introduzione: titolo e testo in nero, spaziatura stretta
\newenvironment{introduzione}{\begingroup\colorlet{csez}{black}\colorlet{crie}{black}%
  \section*{Introduzione}\addcontentsline{toc}{section}{Introduzione}%
  \setlength{\parskip}{3pt plus 1pt}\setlist{itemsep=1pt, topsep=2pt, parsep=0pt}%
  \setlength{\abovedisplayskip}{4pt plus 2pt}\setlength{\belowdisplayskip}{4pt plus 2pt}%
  \setlength{\abovedisplayshortskip}{2pt}\setlength{\belowdisplayshortskip}{3pt}}%
  {\par\endgroup}
% compatibilità con le dispense precedenti
\newenvironment{attenzione}[1][]{\begin{osservazione}[#1]}{\end{osservazione}}

% ---------------------------------------------------------- dimostrazioni a passi
\usepackage[skins, breakable]{tcolorbox}
\newcounter{passo}
\newcommand{\passo}[1]{\par\addvspace{5pt}\stepcounter{passo}%
  \noindent{\bfseries\color{cpasso}Passo \thepasso.}\ \textit{#1}\par\nopagebreak[4]}
\newcommand{\qed}{\ifmmode\qquad\blacksquare\else\unskip\nobreak\ \ensuremath{\blacksquare}\fi}
\newcommand{\iniziobarra}{\setlength{\parskip}{4pt plus 1pt}\setcounter{passo}{0}}
\tcbset{barra/.style={blanker, breakable, parbox=false, left=10pt, right=0pt, top=1pt, bottom=1pt,
  before skip=9pt plus 2pt, after skip=9pt plus 2pt}}
\newtcolorbox{dimostrazione}[1][]{barra, borderline west={1.8pt}{0pt}{cdim},
  before upper={\iniziobarra\textbf{\color{cdim}Dimostrazione\tit{#1}.}\par}}
\newtcolorbox{approfondimento}[1][]{barra, borderline west={1.8pt}{0pt}{cappr},
  before upper={\iniziobarra\textbf{\color{cappr}Approfondimento\tit{#1}.}\ \ignorespaces}}
\newtcolorbox{esempio}[1][]{barra, borderline west={1.8pt}{0pt}{cesm},
  before upper={\iniziobarra\textbf{\color{cesm}Esempio\tit{#1}.}\ \ignorespaces}}
% niente titolo di dimostrazione/approfondimento isolato in fondo alla pagina
\AddToHook{env/dimostrazione/before}{\par\spazio{6\baselineskip}}
\AddToHook{env/approfondimento/before}{\par\spazio{5\baselineskip}}
\AddToHook{env/esempio/before}{\par\spazio{4\baselineskip}}
% conseguenze introdotte dalla domanda a cui rispondono
\newcommand{\domanda}[1]{\par\addvspace{5pt}\spazio{5\baselineskip}\noindent{\color{cris}$\blacktriangleright$}\ \textbf{#1}\par\nopagebreak[4]}
\newcommand{\nsp}[1]{\spazio{#1\baselineskip}}

% ---------------------------------------------------------- riquadri
%  enunciato: riquadro blu con l'enunciato sopra e, dopo \tcblower, il significato
%  \begin{teorema}[di Bloch] enunciato \tcblower significato \end{teorema}
%  (anche: modello). Numerazione comune con Definizione, Osservazione, ...
\newtcolorbox{enunciato}[1]{enhanced, unbreakable, parbox=false,
  colback=cenb, colframe=csez, boxrule=0.6pt, arc=2pt,
  colbacktitle=csez, coltitle=white, fonttitle=\bfseries, title={#1},
  left=9pt, right=9pt, top=6pt, bottom=6pt, toptitle=2pt, bottomtitle=2pt,
  middle=5pt, segmentation style={csez!45, dashed, line width=0.5pt},
  before skip=12pt plus 2pt, after skip=12pt plus 2pt,
  before upper={\setlength{\parskip}{5pt}},
  before lower={\setlength{\parskip}{5pt}\textbf{\color{cint}Cosa dice.}\ \ignorespaces}}
\newenvironment{teorema}[1][]{\par\refstepcounter{enun}\begin{enunciato}{Teorema \theenun\tit{#1}}}{\end{enunciato}}
\newenvironment{modello}[1][]{\par\refstepcounter{enun}\begin{enunciato}{Modello \theenun\tit{#1}}}{\end{enunciato}}
%  chiave: SOLO i risultati essenziali (barra del titolo verde, non spezzabile)
\newtcolorbox{chiave}[1][Risultato]{enhanced, unbreakable, parbox=false,
  colback=cbox, colframe=cris, boxrule=0.6pt, arc=2pt,
  colbacktitle=cris, coltitle=white, fonttitle=\bfseries, title={#1},
  left=9pt, right=9pt, top=5pt, bottom=5pt, toptitle=2pt, bottomtitle=2pt,
  before skip=11pt plus 2pt, after skip=11pt plus 2pt,
  before upper={\setlength{\parskip}{5pt}}}
%  formula da evidenziare senza titolo
\newcommand{\formula}[1]{\par\addvspace{4pt}\begin{center}\begin{tcolorbox}[enhanced, hbox, colback=cform,
  colframe=cformb, boxrule=0.5pt, arc=2pt, left=10pt, right=10pt, top=2pt, bottom=2pt]
  $\displaystyle #1$\end{tcolorbox}\end{center}\par\addvspace{2pt}}
%  informazioni sul corso o curiosità sull'hardware (rarissima)
\newcommand{\titinfo}[1]{\if\relax\detokenize{#1}\relax\else\textbf{#1.}\ \fi}
\newtcolorbox{infobox}[1][]{enhanced, breakable, parbox=false, colback=black!4, frame hidden,
  boxrule=0pt, arc=2pt, left=9pt, right=9pt, top=5pt, bottom=5pt,
  before upper={\setlength{\parskip}{4pt}\titinfo{#1}\ignorespaces}}
%  formulario: riquadro verde spezzabile, una \voce per riga (si spezza tra le voci)
\newtcolorbox{formulario}[1][Formule della lezione]{enhanced, breakable, parbox=false,
  colback=cbox, colframe=cris, boxrule=0.6pt, arc=2pt,
  colbacktitle=cris, coltitle=white, fonttitle=\bfseries, title={#1},
  left=9pt, right=9pt, top=5pt, bottom=7pt, before upper={\setlength{\parskip}{0pt}}}
\newcommand{\gruppo}[1]{\par\addvspace{6pt}\noindent\textbf{#1}\par\nopagebreak\addvspace{4pt}}
\newcommand{\sezvoce}[1]{\if\relax\detokenize{#1}\relax\else\ {\footnotesize\color{black!55}(\S#1)}\fi}
\newcommand{\voce}[3][]{\par\noindent\begin{tabularx}{\linewidth}{@{}>{\raggedright\arraybackslash}p{3.9cm}@{\hspace{8pt}}X@{}}
  #2\sezvoce{#1} & $\displaystyle #3$\end{tabularx}\par\addvspace{4pt}}
%  domande di verifica: \begin{domande}\item \textbf{Domanda?} traccia...\end{domande}
\newlist{domande}{enumerate}{1}
\setlist[domande]{label=\textbf{Q\arabic*.}, leftmargin=2.3em, itemsep=6pt}

% ---------------------------------------------------------- macro generiche
\newcommand{\ingl}[1]{\textit{\color{black!55}(#1)}}   % termine inglese delle slide
\newcommand{\va}{\mathbf a}
\newcommand{\vb}{\mathbf b}
\newcommand{\vk}{\mathbf k}
\newcommand{\vr}{\mathbf r}
\newcommand{\vR}{\mathbf R}
\newcommand{\vG}{\mathbf G}
\newcommand{\vc}{\mathbf c}
\newcommand{\vT}{\mathbf T}
\newcommand{\xh}{\hat{\mathbf x}}
\newcommand{\yh}{\hat{\mathbf y}}
\newcommand{\zh}{\hat{\mathbf z}}
\newcommand{\coord}[1]{\textsf{#1}}
\newcommand{\h}{\tfrac12}
\newcommand{\jj}{\mathrm j}
\newcommand{\ee}{\mathrm e}
\newcommand{\dd}{\mathrm d}
\newcommand{\pd}[2]{\frac{\partial #1}{\partial #2}}
\newcommand{\figura}[1]{\input{figs/#1.tex}}

% ---------------------------------------------------------- float
\renewcommand{\topfraction}{0.9}
\renewcommand{\bottomfraction}{0.8}
\renewcommand{\textfraction}{0.1}
\renewcommand{\floatpagefraction}{0.8}
\setcounter{topnumber}{3}
\setcounter{bottomnumber}{2}
\setcounter{totalnumber}{4}

% ---------------------------------------------------------- aggiunte della Parte 3
\newcommand{\EF}{E_{\mathrm F}}
\newcommand{\kF}{k_{\mathrm F}}
\newcommand{\kB}{k_{\mathrm B}}
\newcommand{\Eg}{E_{\mathrm g}}
\newcommand{\vion}{v_{\mathrm{ion}}}
\newcommand{\Kt}{\tilde K}
\definecolor{occ}{HTML}{2E6FB7}
\definecolor{hole}{HTML}{C0392B}
\definecolor{exc}{HTML}{23884A}
\definecolor{bandfull}{HTML}{8C2A1F}
\definecolor{bandempty}{HTML}{F2E6A0}
\definecolor{bandpart}{HTML}{D9C84A}
%  le tre approssimazioni del §6: stesso riquadro blu di teoremi e modelli
\newenvironment{approssimazione}[1][]{\par\refstepcounter{enun}\begin{enunciato}{Approssimazione \theenun\tit{#1}}}{\end{enunciato}}
\makeatletter
\setlength{\@fptop}{0pt}
\makeatother
